% titel: Blatt 6, A5: Eulersche Formel
% skript: Kap. 1 (Eulersche Formel, komplexe Exponentialfunktion)

% Teil I: Werkzeuge, kein Lösungsweg.

Mit der Eulerschen Formel $e^{it}=\cos t+i\sin t$ wichtige
Eigenschaften der komplexen Exponentialfunktion herleiten und daraus die
Doppelwinkelformeln ableiten.

\begin{werkzeugkasten}[Eulersche Formel und Betrag]
Für alle $t\in\mathbb R$:
\begin{equation*}
  e^{it}=\cos t+i\sin t,\qquad |e^{it}|=\sqrt{\cos^2t+\sin^2t}=1 .
\end{equation*}
Der Zeiger $e^{it}$ läuft also auf dem Einheitskreis. Für das Konjugierte
gilt
\begin{equation*}
  \overline{e^{it}}=\cos t-i\sin t=e^{-it},
\end{equation*}
also ist $e^{-it}$ das Spiegelbild von $e^{it}$ an der reellen Achse.

\textit{Skript:} Kap. 1 (Eulersche Formel).
\end{werkzeugkasten}

\begin{werkzeugkasten}[Induktion und Real-/Imaginärteil]
\begin{itemize}
  \item $(e^{it})^n=e^{int}$ folgt per vollständiger Induktion aus den
        Rechenregeln $e^{ia}e^{ib}=e^{i(a+b)}$ der Exponentialfunktion.
  \item Aus $(e^{it})^2=e^{i2t}$ erhält man durch Ausmultiplizieren der
        linken Seite eine Identität; Vergleich von Real- und Imaginärteil
        liefert die Doppelwinkelformeln.
\end{itemize}
\textit{Skript:} Kap. 1, Satz 3.1.4.
\end{werkzeugkasten}

\begin{achtung}[Darauf achten]
\begin{itemize}
  \item $|e^{it}|=1$ gilt für \emph{reelle} $t$; für komplexe Argumente
        gilt das nicht mehr so einfach (dort wächst das Argumentmodul).
  \item Beim Spiegeln: Konjugation entspricht dem Ersetzen $i\to-i$,
        im reellen Ausdruck ergibt sich $\cos t-i\sin t=e^{-it}$.
  \item Beim Ausmultiplizieren von $(\cos t+i\sin t)^2$ beide Terme
        (Realteil aus $i^2=-1$) korrekt sammeln.
  \item Nur gleiche Frequenzen kombinieren; $n=2$ ergibt genau
        $\cos2t$ und $\sin2t$.
\end{itemize}
\end{achtung}
